\documentclass[a4paper,10pt]{ctexart}
\usepackage[T1]{fontenc}
\usepackage{geometry}
\usepackage{titlesec}
\usepackage{graphicx}
\usepackage{float}
\usepackage{xcolor}

% 页面边距设置
\geometry{left=2cm,right=2cm,top=1.5cm,bottom=1.5cm}

% 移除页眉页脚
\pagestyle{empty}

% 自定义章节格式
\titleformat{\section}{\large\bfseries}{}{0em}{}[\titlerule]
\titlespacing{\section}{0pt}{10pt}{5pt}

% 自定义子章节格式
\titleformat{\subsection}{\bfseries}{}{0em}{}
\titlespacing{\subsection}{0pt}{5pt}{2pt}

% 自定义紧凑列表环境
\newenvironment{tightitemize}
{ \begin{itemize}
    \setlength{\itemsep}{0pt}
    \setlength{\parskip}{0pt}
    \setlength{\parsep}{0pt}
}
{ \end{itemize} }

\begin{document}

\begin{center}
    {\LARGE \textbf{高性能异步多模态智能体架构 (xy-core)}} \\
    \vspace{0.2cm}
    {\large \textit{执行摘要}}
\end{center}

\vspace{0.1cm}
\noindent\textbf{说明}: 本摘要以仓库内 \texttt{legacy/mvp\_core} 的可复现实验基准与实现为准。

\vspace{0.3cm}

\section{一、核心架构：Two-Tier Hybrid（控制平面 vs 数据平面）}
针对 Python 在多模态负载（LLM 推理、VL 视觉理解、图像生成等）中易出现的 GIL 争用与事件循环饥饿（event loop starvation）问题，xy-core 采用 \textbf{Two-Tier Hybrid Architecture}：将编排与路由（控制平面）与计算执行（数据平面）解耦，使控制路径保持轻量、可预测。

\begin{tightitemize}
    \item \textbf{Control Plane (L1 Router)}：FastAPI WebSocket 接入 + \texttt{GlobalTaskScheduler}。负责优先级/队列化、A/B 路由、故障回退（fallback），并对重型任务（例如图像生成）通过队列容量与并发上限实现背压。
    \item \textbf{Data Plane (L2 Compute)}：承载计算密集型负载。默认使用 Python 线程池卸载（thread offloading），并可选启用 C++ 调度器（\texttt{CPPSchedulerEngine/pybind11}）在 GIL 临界路径之外执行与隔离资源。
\end{tightitemize}

\section{二、实测性能数据（RTX 5070 Laptop GPU 8GB，2025-12-23）}
数据来源：\\
\texttt{legacy/mvp\_core/experiment\_results/xy\_core\_real.json} \\
与 \texttt{legacy/mvp\_core/experiment\_results/xy\_core\_cpp\_sched\_real.json}。

\subsection{1. 控制平面响应性（Event Loop Lag）}
\begin{tightitemize}
    \item \textbf{xy-core（Python）}：最大漂移 \textbf{15.72ms}，平均 \textbf{9.55ms}。
    \item \textbf{xy-core（C++ Scheduler）}：最大漂移 \textbf{26.28ms}，平均 \textbf{9.62ms}。
    \item \textbf{结论}：在真实多模态链路下，控制平面仍保持毫秒级漂移区间，未出现长时间阻塞尖峰。
\end{tightitemize}

\subsection{2. 吞吐与端到端时延}
\begin{tightitemize}
    \item \textbf{吞吐（并发 10）}：Python \textbf{0.0386 RPS}；C++ Scheduler \textbf{0.0403 RPS}。
    \item \textbf{P50/P95（Exp3）}：Python \textbf{35.65s / 40.69s}；C++ Scheduler \textbf{30.93s / 34.05s}（中位数约 \textbf{-13\%}，P95 约 \textbf{-16\%}）。
    \item \textbf{结论}：Two-Tier 将重负载从控制平面剥离后，能在不牺牲交互性的前提下提升链路可预测性与尾延迟表现。
\end{tightitemize}

\subsection{3. 可观测性产物 (Observability Artifacts)}
\begin{tightitemize}
    \item \textbf{Trace Timeline}：记录任务排队（scheduled$\to$started）与执行（started$\to$completed）分段，用于定位背压、并行重叠与瓶颈。
    \item \textbf{VRAM Time Series}：通过 NVIDIA-SMI 采样显存占用随时间变化，辅助验证 8GB 显存下的资源隔离策略。
\end{tightitemize}

\section{三、关键技术点}
\begin{tightitemize}
    \item \textbf{准入控制与背压}：对 GPU 重型任务设置队列容量与并发上限，避免无限堆积导致的 OOM 与尾延迟放大。
    \item \textbf{异步卸载与隔离}：对 CPU-bound/GPU-bound 任务分层调度，确保控制平面不承担计算密集型执行。
    \item \textbf{可复现实验链路}：统一输出 Real/Mock 的结果 JSON、图表与 Trace/VRAM 产物，支撑报告级对比。
\end{tightitemize}

\section{四、总结}
xy-core 通过 Two-Tier Hybrid（Control Plane vs Data Plane）将编排与计算解耦，在 8GB 显存的消费级 GPU 上实现了可预测的交互响应性与可解释的执行轨迹。实测数据显示：在真实多模态链路下，控制平面保持毫秒级漂移，同时端到端中位数与尾延迟显著下降，具备工程落地与持续优化的可行性。

\end{document}
